\documentclass[11pt]{article}
\usepackage{fontspec}
\usepackage[english,ngerman]{babel}
\usepackage[margin=25mm]{geometry}
\usepackage{hyperref}
\begin{document}
\raggedright
\selectlanguage{english}
\section*{Translation sources and terminology}
These English and German texts are translated from Gauss's original-language editions. Earlier translations are consulted to standardize mathematical terminology and usage; they are not reproduced as this edition's running text. The original mathematical formulations remain those of Gauss's source.
\begin{itemize}
\item H. Maser, German editor and translator, \emph{Carl Friedrich Gauss' Untersuchungen über höhere Arithmetik}, Berlin: Julius Springer, 1889. Reference for German arithmetical terminology and corresponding passages. \url{https://archive.org/details/carlfriedrichga00gausgoog}
\item James Caddall Morehead and Adam Miller Hiltebeitel, translators, \emph{General Investigations of Curved Surfaces of 1827 and 1825}, Princeton University Library, 1902; Project Gutenberg 36856. Reference for English mathematical terminology and usage. \url{https://www.gutenberg.org/ebooks/36856}
\item Paul Taylor with Bernard Leak, English translation of Gauss's second proof of the fundamental theorem of algebra, 1983. Reference for English mathematical terminology and usage. \url{https://www.paultaylor.eu/misc/gauss.html}
\item Ulrich Görtz, \emph{Algebra}, Wintersemester 2021/22, Section 3.1; \emph{Lineare Algebra II}, Sommersemester 2021, Section 18.4. References for contemporary German terminology. \url{https://math.ug/a-ws2122/sect0006.html} and \url{https://math.ug/la2-ss21/sec-quotient-ring.html}
\item Nikolay Moshchevitin, \emph{Eine Bemerkung über positiv definite quadratische Formen und rationale Punkte}, arXiv:1510.06214v3. Reference for German quadratic-form terminology. \url{https://arxiv.org/abs/1510.06214v3}
\end{itemize}
The geometrical and algebraic references are used for language in their respective subjects. They do not supply replacement text or proofs for these arithmetic papers.
\selectlanguage{ngerman}
\section*{Quellen und Terminologie der Übersetzungen}
Die englischen und deutschen Texte sind aus Gauss' originalsprachlichen Ausgaben übersetzt. Frühere Übersetzungen dienen dem Abgleich der mathematischen Terminologie und des Sprachgebrauchs; ihr Wortlaut wird nicht als fortlaufender Text dieser Ausgabe übernommen. Die mathematischen Formulierungen bleiben die der Originalvorlage. Die oben genannten Übersetzer und Autoren werden für diese terminologische Nutzung gewürdigt.
\end{document}
